\section{Source bounds and retained routes}
\label{sec:global-construction}

We apply the construction of Zhao--Booth-Clibborn--Schwahn
\cite[Sections~4--5]{ZBS}, based on Lu--Wang's retained-set descent
\cite[Theorem~2]{LW}. A \emph{source} is a positive pair $(A,B)$ for
which $R(k,\ell)\le Ce^{Ak+B\ell}$ holds with one $C\ge1$ for all
positive $k,\ell$. A \emph{profile} $P=(p_P,I_P,L_P)$ has
$0<p_P<1$, a compact interval $I_P\subset(0,\infty)$, and a positive
height function $L_P$. With $\rho=\ell/k$ and $d_R(G)$ denoting red
edge density, its assertion is
\begin{equation}\label{eq:profile}
 d_R(G)\ge p_P,\quad \rho\in I_P,\quad
 |V(G)|\ge Ce^{kL_P(\rho)}
 \quad\Longrightarrow\quad
 G\text{ contains a red }K_k\text{ or a blue }K_\ell,
\end{equation}
with $C$ uniform in the targets and graph.

A \emph{control} $\gamma$ consists of parameters satisfying
Proposition~\ref{prop:local-alternative}, a proposed terminal source
$(A_\sigma,B_\sigma)$, and
$0<\xi<\min\{x,e^{-A_\sigma}\}$, $0<y_0<e^{-B_\sigma}$,
$0<\nu<\mu$. Put $a_\gamma=-\log\xi$, $b_\gamma=-\log y_0$ and
$c_\gamma=-\log\nu$.
The proof of~\cite[Theorem~4.2, Section~4.3]{ZBS} applies with our local
alternative substituted. It uses the terminal source only on proper
subsets of the current host. We use this closure result and the leaf
and route construction of~\cite[Definition~5.4 and Lemma~5.8]{ZBS};
the closure induction originates in~\cite[Lemma~A.1]{LW}.

A \emph{retained route} has knots
$0<\theta=u_0<\cdots<u_M=\omega$ and a continuous piecewise-affine
function $g$ with $g(\theta)=g_0>0$ and slope $d_j$ on
$[u_j,u_{j+1}]$. Choose a child profile $P_j$ whose domain contains
that interval, with $L_j<g$ there and $d_j>-\log(1-p_j)$.
Writing $w=w_\gamma$ for the control's weight, the leaf and route height bounds are
\begin{equation}\label{eq:route-output}
 \begin{aligned}
 L_\gamma(\rho)&=\frac{a_\gamma+\rho(b_\gamma+wc_\gamma)}{1+w},\\
 V_\gamma(\rho)&=\max\left\{g(\rho),\,
 \frac{a_\gamma+\rho b_\gamma+
 w\bigl(\theta c_\gamma+g(\rho)-g_0\bigr)}{1+w}\right\}.
 \end{aligned}
\end{equation}
To establish a profile $P$ on an interval $J$, require
$L_P\ge L_\gamma+\epsilon$ for a leaf or
$L_P\ge V_\gamma+\epsilon$ for a route throughout $J$, for some
$\epsilon>0$. Also require $p_P\ge p_\gamma$ for a leaf and
$p_P>p_\gamma$ for a route, with $J\subseteq[\theta,\omega]$ in the latter case.
Each child profile is applied on a proper subhost.

Take finite families of proposed sources and profiles. Profiles built
from leaves, routes and proper-subhost complementary pairs are
\emph{direct}; those deduced from direct profiles by density transfer,
complementary pairs and monotonicity in the blue target are \emph{combined}.
Proper-subhost complementary pairs and target-monotonicity extend the rules
of~\cite[Definitions~5.4--5.5]{ZBS}.
Every profile domain must have a finite cover by valid pieces.
For each proposed source $(A,B)$, require some $\epsilon>0$ such that,
at every $\rho>0$, either the symmetric Lu--Wang baseline $F$ or a
pair of profiles with $p_R+p_B\le1$ satisfies
\[
 F(\rho)\le A+B\rho-\epsilon(1+\rho)
 \quad\text{or}\quad
 \max\{L_R(\rho),\rho L_B(1/\rho)\}
 \le A+B\rho-\epsilon(1+\rho).
\]
The profiles must be defined at the indicated ratios. Lu--Wang's
baseline is uniform over all positive target pairs, with arbitrarily
small positive error proportional to $k+\ell$~\cite[Lemma~B.1]{LW}.

\begin{proposition}[Joint construction]\label{prop:joint-construction}
For a finite family satisfying the profile conditions and source comparisons
above, there is a common constant $C\ge1$ for which every source bound and
profile assertion holds.
\end{proposition}
\begin{proof}
Apply the induction of~\cite[Theorem~5.7]{ZBS} with the local alternative
of Proposition~\ref{prop:local-alternative}.
Choose a common prefactor for the small-target cases
and baseline estimates. At host order $N$, the induction hypothesis
supplies every assertion on proper subhosts; prove the direct profiles,
then the combined profiles, then the sources at order $N$.
Strong induction on $N$ proves all the assertions.
\end{proof}
